\problem{1}
$\bigl[\textbf{Омск, 2024, 9.3}\bigr]$
Какой цифрой оканчивается число $2024^n+2024^{n+3}$, где $n$ - натуральное число? Ответ поясните.

\problem{2}
$\bigl[\textbf{ДНР 2024 9.6}\bigr]$
Простое число $p$ таково, что числа $15 p-15$ и $16 p-15$ являются точными квадратами. Чему может быть равно $p$ ? Укажите все возможные варианты в любом порядке.

\problem{3}
$\bigl[\textbf{Москва 2024 10.4.}\bigr]$ Есть 30 нечётных натуральных чисел. Оказалось, что их можно разбить на 10 троек так, что в каждой тройке одно из чисел равно произведению двух других. Может ли сумма этих 30 чисел равняться 2024?

\problem{4}
$\bigl[\textbf{Питер 2021 11.2}\bigr]$
Сергей выписал числа от 500 до 1499 в строчку в некотором порядке. Под каждым числом, кроме самого левого, он написал НОД этого числа и его левого соседа, получив вторую строчку из 999 чисел. Далее он по тому же правилу получил из неё третью строчку из 998 чисел, из неё - четвертую из 997 чисел и и. д. Остановился он, когда впервые все числа в очередной строчке оказались равны единице. Какое наибольшее количество строчек он мог выписать к этому моменту?

\problem{5}
$\bigl[\textbf{Лен область 2022, 11.5}\bigr]$
Компьютер сгенерировал несколько различных натуральных чисел. Для каждой пары чисел он выяснил, на какую наибольшую степень числа 2022 делится их разность. Вариант, что разность делится только на нулевую степень, тоже возможен. Оказалось, что разных ответов у компьютера получилось 2023. Какое наибольшее количество чисел мог сгенерировать компьютер?